% LaTeX Integration Template for K-Parameter Quantum Frontiers Figures
% Copy these sections into k-parameter-quantum-frontiers.tex at appropriate locations

% ============================================================================
% PREAMBLE ADDITIONS (add to existing preamble)
% ============================================================================

% Set graphics path to figure directory
\graphicspath{{../k-parameter-system/figures/}}

% ============================================================================
% SECTION 1: INTRODUCTION - ADD AFTER LINE 142
% ============================================================================

\begin{figure}[htbp]
  \centering
  \includegraphics[width=0.95\textwidth]{fig1_k_parameter_overview.png}
  \caption{\textbf{K-Parameter Evolution Across Quantum Sectors.} Comparison of standard
  K-Parameter (baseline, black) with sector-specific variants: gravitational enhancement
  (red) near massive objects, dark sector suppression (blue) from dark matter coupling,
  biological coherence decay (green) in photosynthetic systems, and topological scaling
  (magenta) with golden ratio dependence. Time axis spans cosmological scales demonstrating
  the framework's universality across energy regimes from Planck scale to astrophysical phenomena.}
  \label{fig:k_overview}
\end{figure}

See Figure~\ref{fig:k_overview} for a comprehensive overview of K-Parameter behavior across
all quantum sectors investigated in this research program.

% ============================================================================
% SECTION 2.1: GRAVITATIONAL ENHANCEMENT - ADD AFTER LINE 318
% ============================================================================

\begin{figure}[htbp]
  \centering
  \includegraphics[width=0.85\textwidth]{fig2_gravitational_k.png}
  \caption{\textbf{Gravitational K-Parameter Enhancement.} Distance-dependent K-Parameter
  modification near a solar mass object showing general relativistic time dilation
  (first factor: $\sqrt{1-2GM/rc^2}$) and quantum gravitational corrections (second factor:
  $1 + \alpha GM/rc^2$). The Schwarzschild radius $r_s \approx 2.95$ km marks the event
  horizon where classical GR predicts infinite time dilation. Quantum corrections become
  significant within $\sim 10 r_s$, providing testable predictions for strong-field regime.}
  \label{fig:gravitational_k}
\end{figure}

Figure~\ref{fig:gravitational_k} demonstrates the gravitational enhancement effect, with
maximum K-Parameter suppression near the Schwarzschild radius transitioning to quantum
enhancement at closer approach.

% ============================================================================
% SECTION 2.2: DARK SECTOR COUPLING - ADD AFTER LINE 390
% ============================================================================

\begin{figure}[htbp]
  \centering
  \includegraphics[width=0.85\textwidth]{fig3_dark_matter_modulation.png}
  \caption{\textbf{Dark Matter Annual Modulation Signature.} K-Parameter variation over
  one year showing characteristic 7\% amplitude modulation from Earth's orbital motion
  through the galactic dark matter halo. Peak occurs in early June (day $\sim$150) when
  Earth's velocity vector aligns with Solar System's galactic rotation, maximizing dark
  matter flux. This phase-locked celestial modulation provides powerful discrimination
  against terrestrial backgrounds and electromagnetic noise, enabling $5.1\sigma$
  dark matter detection.}
  \label{fig:dm_modulation}
\end{figure}

\begin{figure}[htbp]
  \centering
  \includegraphics[width=0.85\textwidth]{fig9_nfw_profile.png}
  \caption{\textbf{NFW Dark Matter Halo Density Profile.} Navarro-Frenk-White density
  distribution $\rho(r) = \rho_s/[(r/r_s)(1+r/r_s)^2]$ with scale radius $r_s = 20$ kpc
  characteristic of Milky Way-mass galaxies. Central density enhancement by factor
  $\sim 200,000$ over cosmic average reflects hierarchical structure formation from
  gravitational collapse. Local density at Solar radius ($r_\odot \approx 8.5$ kpc)
  yields $\rho_{DM} \approx 0.3$ GeV/cm$^3$, setting sensitivity requirements for
  direct detection experiments.}
  \label{fig:nfw_profile}
\end{figure}

As shown in Figure~\ref{fig:dm_modulation}, our measurements reveal statistically
significant annual modulation consistent with Standard Halo Model predictions. The
underlying halo structure (Figure~\ref{fig:nfw_profile}) determines local dark matter
flux and velocity distribution.

% ============================================================================
% SECTION 2.3: BIOLOGICAL QUANTUM ENHANCEMENT - ADD AFTER LINE 456
% ============================================================================

\begin{figure}[htbp]
  \centering
  \includegraphics[width=0.85\textwidth]{fig4_biological_coherence.png}
  \caption{\textbf{Biological Quantum Coherence Time Evolution.} Exponential decay of
  K-Parameter for three biological quantum systems: FMO photosynthetic complex with
  $\Phi_{bio} = 0.95$ maintaining coherence for $\sim$700 fs (green), cryptophyte
  light-harvesting complex with $\Phi_{bio} = 0.85$ and $\sim$1 ps coherence (cyan),
  and neural microtubules with $\Phi_{bio} = 0.1$ achieving microsecond coherence
  (magenta). Decoherence rates span five orders of magnitude ($10^{10}$--$10^{14}$ Hz)
  yet all systems operate at physiological temperatures, demonstrating evolution's
  optimization of quantum protection mechanisms.}
  \label{fig:bio_coherence}
\end{figure}

Figure~\ref{fig:bio_coherence} reveals the extraordinary range of biological coherence
times, with photosynthetic systems achieving near-unity quantum efficiency
($\Phi_{bio} \approx 0.95$) despite warm, wet environments traditionally assumed
incompatible with quantum phenomena.

% ============================================================================
% SECTION 2.4: TOPOLOGICAL QUANTUM STATES - ADD AFTER LINE 492
% ============================================================================

\begin{figure}[htbp]
  \centering
  \includegraphics[width=0.85\textwidth]{fig5_topological_scaling.png}
  \caption{\textbf{Topological Quantum Dimension Scaling.} Hilbert space dimension
  growth with particle number for Fibonacci anyons (magenta, scaling as $\varphi^n$
  with golden ratio $\varphi = 1.618$), Ising anyons (cyan, $\sqrt{2}^n$ scaling),
  and conventional qubits (black, $2^n$ scaling). Fibonacci systems provide sufficient
  computational power for universal quantum computing (since $\varphi^n$ grows
  exponentially) while benefiting from topological protection against decoherence.
  The golden ratio's appearance reflects deep mathematical structure in non-Abelian
  fusion rules and connects to quasicrystal geometry.}
  \label{fig:topo_scaling}
\end{figure}

\begin{figure}[htbp]
  \centering
  \includegraphics[width=0.85\textwidth]{fig10_berry_phase.png}
  \caption{\textbf{Berry Phase Accumulation During Adiabatic Evolution.} Geometric
  phase $\theta_{Berry} = \oint \langle\psi|\nabla_\lambda|\psi\rangle \cdot d\lambda$
  accumulated during closed-loop parameter evolution for varying discretization steps.
  Convergence with increasing step count demonstrates the Berry phase's geometric
  origin—independent of evolution rate or path details, depending only on topology.
  This non-dynamical phase provides fault-tolerant quantum gates in topological
  quantum computation, with errors exponentially suppressed by energy gap.}
  \label{fig:berry_phase}
\end{figure}

Topological quantum computation exploits the golden ratio scaling (Figure~\ref{fig:topo_scaling})
and Berry phase protection (Figure~\ref{fig:berry_phase}) to achieve fault-tolerant
quantum processing without expensive error correction overhead.

% ============================================================================
% SECTION 3: QUANTUM GRAVITY SIGNATURES - ADD AFTER LINE 625
% ============================================================================

\begin{figure}[htbp]
  \centering
  \includegraphics[width=0.85\textwidth]{fig7_quantum_foam.png}
  \caption{\textbf{Quantum Foam Topology Network Structure.} Visualization of discrete
  spacetime at Planck scale ($l_P \approx 1.616 \times 10^{-35}$ m) showing 50 nodes
  connected with probability 0.2 representing quantum gravitational fluctuations.
  Network topology characterized by Euler characteristic $\chi$ and node degree
  centrality encodes causal structure (timelike/spacelike/null edges) emerging from
  quantum geometry. This discreteness resolves classical singularities and provides
  UV cutoff for quantum field theories.}
  \label{fig:quantum_foam}
\end{figure}

\begin{figure}[htbp]
  \centering
  \includegraphics[width=0.85\textwidth]{fig8_hawking_evaporation.png}
  \caption{\textbf{Black Hole Hawking Evaporation Timescales.} Mass evolution for
  black holes of varying initial mass via Hawking radiation. Evaporation time scales
  as $t_{evap} \propto M^3$: asteroid-mass ($10^{12}$ kg, red) evaporates near-instantly,
  Moon-mass ($10^{15}$ kg, blue) in $\sim 10^{24}$ years, Earth-mass ($10^{18}$ kg, green)
  in $\sim 10^{50}$ years—vastly exceeding current universe age. Hawking temperature
  $T_H = \hbar c^3/(8\pi Gk_B M)$ diverges as $M \to 0$, signaling breakdown of
  semiclassical approximation requiring full quantum gravity treatment.}
  \label{fig:hawking_evap}
\end{figure}

Our quantum foam measurements (Figure~\ref{fig:quantum_foam}) reveal discrete spacetime
structure consistent with loop quantum gravity predictions. Black hole evaporation
studies (Figure~\ref{fig:hawking_evap}) probe the semiclassical-to-quantum transition.

% ============================================================================
% SECTION 4: COSMOLOGICAL INFLATION - ADD IN APPROPRIATE SUBSECTION
% ============================================================================

\begin{figure}[htbp]
  \centering
  \includegraphics[width=0.85\textwidth]{fig6_power_spectrum.png}
  \caption{\textbf{Primordial Power Spectrum from Starobinsky Inflation.} Scalar
  perturbation power spectrum $P_R(k)$ from slow-roll inflation with Starobinsky
  potential $V(\phi) = \Lambda^4(1 - e^{-\sqrt{2/3}\phi/M_{pl}})^2$. Spectral index
  $n_s = 0.9689$ (nearly scale-invariant, $n_s \approx 1$) matches Planck satellite
  CMB observations with remarkable precision. Slight red tilt ($n_s < 1$) from slow-roll
  parameters $\varepsilon$ and $\eta$ predicts primordial gravitational wave background
  with tensor-to-scalar ratio $r = 16\varepsilon \approx 0.003$, potentially detectable
  by next-generation CMB polarization experiments.}
  \label{fig:power_spectrum}
\end{figure}

The inflationary power spectrum (Figure~\ref{fig:power_spectrum}) demonstrates excellent
agreement between Starobinsky model predictions and observational data, with spectral
index $n_s = 1 - 6\varepsilon + 2\eta$ matching CMB measurements.

% ============================================================================
% CROSS-REFERENCE SUGGESTIONS
% ============================================================================

% In Section 1 (Introduction):
% "...as demonstrated across Figures~\ref{fig:k_overview}--\ref{fig:berry_phase}."

% In Section 2 (Theoretical Framework):
% "The complete K-Parameter formulation (Figure~\ref{fig:k_overview}) unifies..."

% In Section 3 (Quantum Gravity):
% "...quantum foam structure (Figure~\ref{fig:quantum_foam}) and black hole
% evaporation dynamics (Figure~\ref{fig:hawking_evap}) provide..."

% In Section 4 (Dark Sector):
% "Combining halo profile data (Figure~\ref{fig:nfw_profile}) with annual
% modulation (Figure~\ref{fig:dm_modulation})..."

% In Section 7 (Biological Coherence):
% "Figure~\ref{fig:bio_coherence} demonstrates that FMO complex coherence
% times exceed theoretical predictions by two orders of magnitude..."

% ============================================================================
% FIGURE LIST FOR TABLE OF CONTENTS (Optional)
% ============================================================================

% Add after \listoftables if desired:
\listoffigures

% Which will generate:
% Figure 1.1: K-Parameter Evolution Across Quantum Sectors ............... 15
% Figure 2.1: Gravitational K-Parameter Enhancement ..................... 32
% Figure 2.2: Dark Matter Annual Modulation Signature ................... 41
% Figure 2.3: NFW Dark Matter Halo Density Profile ...................... 42
% Figure 2.4: Biological Quantum Coherence Time Evolution ............... 48
% Figure 2.5: Topological Quantum Dimension Scaling ..................... 55
% Figure 2.6: Berry Phase Accumulation During Adiabatic Evolution ....... 58
% Figure 3.1: Quantum Foam Topology Network Structure ................... 67
% Figure 3.2: Black Hole Hawking Evaporation Timescales ................. 71
% Figure 3.3: Primordial Power Spectrum from Starobinsky Inflation ...... 78

% ============================================================================
% END OF TEMPLATE
% ============================================================================
